\chapter{How to answer anything}
\label{ch:answering}

\section*{What this chapter is for}

A method for answering a technical question you have not prepared for, which is
most of them. It is four steps, it works on any question in this book's bank
and on questions the bank does not contain, and it is the thing to practise
rather than the answers.

\section{Why a method rather than answers}

A bank of answers has two failure modes. It runs out, at the first question
nobody anticipated. And it produces recall, which sounds like recall \dash{}
fluent, complete, and audibly unconnected to the situation in front of you.

\judgement{The strongest tell that somebody is reciting is that the answer does
not change when the question's constraints change. An interviewer who adds
``and this runs on one machine with fifty users'' and gets the same answer has
learnt that the answer was not about their system.}

What survives is a method that produces an answer from the question's own
constraints.

\section{The four steps}

\begin{kitmethod}
\textbf{1. What does it change for the person using this?}

Start at the consequence, not the mechanism. Not ``a B-tree index orders the
keys'' but ``this query currently scans the table, so it gets slower as the
table grows, and the person waiting on it will notice at about a million rows''.

This step is what makes an answer sound like it came from somebody who has
operated a system.

\textbf{2. What does it cost to reverse?}

Sort every option by how hard it is to undo. Adding an index is reversible in
minutes. Choosing an identifier scheme is reversible after a migration.
Publishing an API is reversible after everybody else's release cycle.

\judgement{This step is where Staff-level judgement actually lives. Decisions
that are cheap to reverse deserve a quick answer and a willingness to be wrong;
decisions that are expensive deserve the time. Treating those two the same way
\dash{} in either direction \dash{} is the most common judgement failure at
this level.}

\textbf{3. What is the boring default?}

Name the obvious, well-understood option first and say why it is the default.
Then say what would make you leave it.

An answer that starts at the interesting option has skipped the argument.
Starting at the boring one and departing from it deliberately \emph{is} the
argument, and it takes one extra sentence.

\textbf{4. What would change my mind?}

Name the measurement, the constraint, or the fact that would move you to the
other option.

This is the strongest sentence in the method and the one people leave out. It
converts a preference into a decision, it shows you know where the boundary is,
and it invites the interviewer to supply the constraint \dash{} which is
frequently what they were waiting to do.
\end{kitmethod}

\section{What it sounds like}

The question: \emph{would you cache this?}

Without the method: ``Yes, probably a distributed cache with a short expiry,
though you have to be careful about invalidation.'' True, generic, could have
been said before the question.

With it: ``What it changes for the user is the tail latency, so first I would
want to know whether the slow requests are the repeated ones \dash{} caching a
long tail of unique requests buys nothing. Caching is cheap to add and
expensive to remove, because things start depending on the latency it provides,
so I would rather be slow and correct first. The boring default is no cache
until the numbers say otherwise, and the number that would change my mind is a
hit rate above roughly half on real traffic. If it is a tenth, the cache is
complexity with no payoff.''

Four steps, one paragraph, and it is about their system.

\section{When you do not know}

The method still applies, with one addition: say where the edge of your
knowledge is, and then keep going from first principles.

``I have not used that specific one. What I would want to know is whether it
guarantees ordering, because the whole design changes at that fork \dash{} if
it does, this becomes straightforward; if it does not, I need an idempotent
consumer and that is a different piece of work.''

\judgement{This is a strong answer, and candidates suppress it because earlier
in their careers it was not. At Junior, not knowing is a gap. At Staff, knowing
which question would resolve it is the skill being assessed, and admitting the
boundary is what makes the rest of your answers credible.}

\section{How the bank is structured}

Every question in the chapters that follow has six parts, in the same order
every time.

\begin{kittable}{@{}lX@{}}
\textbf{As asked} & the question in the words it usually arrives in \\
\textbf{What is being tested} & which is rarely the topic named \\
\textbf{The Staff answer} & not the correct answer \dash{} the one that assesses at this level \\
\textbf{The Senior answer that falls short} & correct, and assessed a level lower \\
\textbf{Where this is contested} & because most of these have a real disagreement in them \\
\textbf{Drill} & a variant to answer out loud \\
\end{kittable}

\judgement{The fourth row is the most useful and the most unusual. Most
preparation material gives the right answer; the difference between levels is
rarely correctness, and seeing the good-but-not-senior version beside the
senior one is the fastest way to hear the gap.}

\begin{drill}
Take the last technical question you were asked in any interview and answer it
again, out loud, in four steps. Time it. If it runs past ninety seconds you are
explaining the mechanism rather than the decision \dash{} the four steps should
compress an answer, not extend it.
\end{drill}

\section*{What this chapter does not claim}

The method is not measured and is not original in its parts \dash{} reversibility
as a decision axis in particular is widely used. What is specific here is the
ordering, and the claim that starting at the consequence rather than the
mechanism is what makes an answer sound operational.

It does not claim the method produces correct answers. It produces answers that
are about the question asked, which is a different and more achievable thing.
